\section{Conclusion and Future Work}
\label{sec:conclusion}

\subsection{Summary}

This project implemented and evaluated two classical power side-channel attacks
---DPA and CPA---targeting AES-128 on an RP2040 microcontroller.

\begin{itemize}
  \item A complete hardware test bench was built: Pico firmware with a hardware
        trigger, pyvisa-based oscilloscope acquisition, and a structured
        \texttt{.npy} dataset format.

  \item CPA with the Pearson correlation distinguisher recovers all 16 bytes
        of the round-10 key at rank 1 using 2{,}000 traces.
        Fifteen bytes succeed under the HW model ($c=0$); byte 11 requires
        an HD model ($c=2$).

  \item DPA with the differential-mean distinguisher recovers all 16 bytes
        at bit position $\ell=1$ using the same 2{,}000 traces.

  \item Both attacks recover the original 128-bit key
        (\texttt{2B 7E 15 16 28 AE D2 A6 AB F7 15 88 09 CF 4F 3C})
        via backward key schedule.

  \item The byte-11 HD leakage finding provides a concrete example of
        how leakage model selection affects CPA but not DPA, highlighting a
        practical advantage of DPA in model-agnostic scenarios.
\end{itemize}

\subsection{Future Work}

\begin{description}
  \item[Guessing Entropy on real hardware traces.]
    Once acquisition on the physical Pico is complete
    (\texttt{attacks/guessing\_entropy.py --data data/set\_001\_pico\_real}),
    GE vs.\ $N_\text{traces}$ curves for both attacks will provide a
    quantitative trace-complexity comparison.

  \item[Masking countermeasure evaluation.]
    Applying Boolean masking (first-order) to the firmware S-Box lookup
    and measuring the increase in required traces for second-order attacks
    would directly address the defense portion of the proposal.

  \item[SNR characterisation across noise levels.]
    Injecting controlled additive noise into simulated traces and measuring
    at what SNR each attack fails would produce the characterisation proposed
    in the original plan.

  \item[EM side channel.]
    Replacing the power probe with an EM probe positioned over the RP2040
    core would allow comparison of power vs.\ EM leakage on the same target.
\end{description}
